\section{Kết luận}

\subsection{Kết luận}

Báo cáo đã trình bày tensor factorization như một công cụ thống nhất cho ba
hướng ứng dụng trong học máy: nén và thích nghi mô hình học sâu, suy luận hiệu
quả với probabilistic circuits, và khai phá dữ liệu đa chiều có khả năng diễn
giải. Phần cơ sở lý thuyết cung cấp các khái niệm chính như CP, Tucker,
matricization, Khatri-Rao product, tensor rank và tính duy nhất của CP; các phần
ứng dụng sau đó dùng lại cùng hệ ký hiệu này để giải thích vì sao phân rã tensor
có ý nghĩa thực tế.

Ở phần Deep Learning, thí nghiệm LoRA rank sweep cho thấy low-rank adaptation có
thể đạt hiệu năng gần full fine-tuning với số tham số cập nhật rất nhỏ. Cụ thể,
LoRA rank 8 đạt validation accuracy $0.7511$ trên SST-2 trong khi chỉ cập nhật
$0.1923\%$ tổng số tham số của mô hình. Ở phần Probabilistic Circuit, demo CPJ
circuit kiểm chứng được hai tính chất quan trọng: mô hình học gần đúng phân phối
toy 4-bit với KL $0.0054$ và vẫn mở rộng được sang density estimation trên
Binarized MNIST với test bpd $0.3704$ trong thiết lập nhỏ. Ở phần Data Mining,
các thực nghiệm CP-ALS, PARAFAC2 và CMTF minh họa cách tensor factorization giữ
lại cấu trúc đa mode, từ đó hỗ trợ diễn giải pattern tốt hơn so với việc ép dữ
liệu về dạng ma trận.

\subsection{Các hướng phát triển trong tương lai}

Các hướng phát triển trong tương lai có thể được tổng hợp từ ba mạch chính của
báo cáo. Trước hết, trong học sâu, vấn đề trung tâm không chỉ là áp dụng một phân
rã cụ thể, mà là thiết kế nguyên tắc chọn mô hình và chọn rank. CP rank,
multilinear rank, Tensor Train rank, LoRA rank, expert rank hay width của circuit
đều điều khiển cùng một đánh đổi: biểu diễn càng gọn thì chi phí càng thấp, nhưng
rank quá nhỏ có thể làm mô hình thiếu năng lực; rank quá lớn lại làm mất lợi thế
nén và tăng nguy cơ overfit. Vì vậy, một hướng quan trọng là phát triển các tiêu
chuẩn chọn rank có cơ sở hơn, kết hợp cả độ chính xác, số tham số, bộ nhớ, thời
gian suy diễn, và độ ổn định khi thay đổi cấu trúc phân rã.

Thứ hai, tensorization trong học sâu có thể mở rộng ra ngoài nén trọng số tĩnh.
Báo cáo đã nêu nhiều đối tượng tự nhiên có cấu trúc tensor: activation, attention
score, gradient, update khi fine-tuning, tập expert trong MoE, hoặc nhiều layer và
head trong Transformer. Các hướng như LoRTA, TensorGRAD, tensorized convolution,
polynomial/self-attention hiệu quả và tensor product attention cho thấy tensor
factorization có thể trở thành một nguyên lý thiết kế kiến trúc, không chỉ là
bước hậu xử lý sau khi mô hình đã được huấn luyện. Câu hỏi mở là xác định khi nào
một đối tượng trong mạng có mode tự nhiên đủ rõ để factor thu được vừa tiết kiệm
vừa có ý nghĩa, thay vì chỉ là một phép reshape nhân tạo khó diễn giải.

Thứ ba, ở phía probabilistic circuits, mối liên hệ giữa circuit và non-negative
tensor factorization mở ra nhiều câu hỏi lý thuyết. Các thuộc tính như
smoothness và decomposability giúp suy luận chính xác trở nên khả giải, nhưng
không tự động bảo đảm mô hình luôn biểu diễn gọn mọi phân phối. Do đó các hướng
phát triển gồm hiểu rõ hơn quan hệ giữa expressiveness, compactness và
tractability; học đồng thời tham số và cấu trúc circuit; nghiên cứu điều kiện
identifiability cho deep circuits; xây dựng các thuật toán approximate inference
dựa trên thuộc tính cấu trúc; và mở rộng circuit lên miền liên tục thông qua
integral circuits hoặc quasi-tensors. Đây là bước nối tự nhiên từ tensor rời rạc
sang mô hình xác suất trên hàm và biến liên tục.

Cuối cùng, trong khai phá dữ liệu, hướng phát triển nổi bật là đưa tensor
factorization vào các bài toán đa nguồn, không đầy đủ và có cấu trúc miền ứng
dụng rõ ràng. CP cung cấp tính duy nhất tương đối cho pattern discovery,
PARAFAC2 xử lý các slice có kích thước hoặc căn chỉnh khác nhau, còn CMTF hỗ trợ
phân tích chung nhiều ma trận/tensor thông qua factor dùng chung. Các bài toán
tương lai cần xử lý tốt hơn missing data, dữ liệu lệch thời gian, prior knowledge
từ chuyên gia, ràng buộc không âm/thưa/mượt, và lựa chọn số component. Điểm then
chốt là giữ được cân bằng giữa interpretability/uniqueness và độ linh hoạt của
mô hình: càng thêm phi tuyến hoặc cấu trúc phức tạp, ta càng cần cơ chế kiểm tra
độ ổn định và ý nghĩa của các factor thu được.

Nhìn chung, các hướng trên cho thấy tensor factorization không chỉ là một công
cụ giảm chiều. Nó có thể đóng vai trò như một ngôn ngữ chung để thiết kế mô hình
gọn hơn, suy luận xác suất hiệu quả hơn, và khai phá các pattern có khả năng
diễn giải trong dữ liệu đa chiều, đa nguồn.

\subsection{Bảng phân công công việc}

\begin{table}[H]
\centering
\renewcommand{\arraystretch}{1.2}
\begin{tabular}{p{0.16\textwidth}p{0.24\textwidth}p{0.50\textwidth}}
\toprule
\textbf{MSSV} & \textbf{Họ tên} & \textbf{Phần đóng góp} \\
\midrule
23120119 & Đỗ Tiến Đạt & Phần Probabilistic Circuit: tìm hiểu cơ sở lý thuyết chung của Tensor và Probabilistic Circuit, triển khai thực nghiệm, phân tích kết quả và hỗ trợ tổng hợp báo cáo. \\
23120137 & Triệu Tuấn Kiệt & Phần Deep Learning: tìm hiểu LoRA/LoRTA, triển khai thí nghiệm, phân tích LoRA rank sweep và trực quan hóa kết quả và hỗ trợ viết tổng hợp báo cáo. \\
23120163 & Lê Nhật Minh Tâm & Phần Deep Learning: xử lý thiết lập thực nghiệm, đánh giá kết quả fine-tuning, so sánh với cơ sở lý thuyết và hỗ trợ viết tổng hợp báo cáo. \\
23120134 & Nguyễn Đăng Khoa & Phần Data Mining, giải thích cơ sở lý thuyết, xây dựng thuật toán CP-ALS và CP=OPT, phân tích kết quả và tổng hợp báo cáo. \\
23120105 & Huỳnh Mạnh Tường & Phần Data Mining, xây dụng mô hình PARAFAC2 và CMTF, trình bày ứng dụng thực tiễn, hỗ trợ trình bày báo cáo và thực nghiệm \\
\bottomrule
\end{tabular}
\caption{Bảng phân công công việc của nhóm.}
\label{tab:work-allocation}
\end{table}
